% GENERATED FILE. Edit canonical lessons, then run npm run generate:books.
% canonical-source-hash: fnv1a64:dad5595b4a34d8f3
% canonical-lessons: ML-C162-karati, ML-C162-kurukkan, ML-C162-ottakam, ML-C162-taravu, ML-C162-munna

\chapter{Bear, Fox, Camel, Duck, Owl}
\label{ch:ml-bear-fox-camel-duck-owl}
\hlchaptermodality{162}

\begin{chapteropening}
\textbf{By the end of this chapter:} \emph{I can say a bear, a fox, a camel, a duck and an owl.}

Complete the last lesson of chapter 162: I can say a bear, a fox, a camel, a duck and an owl.
\end{chapteropening}

\section[\texorpdfstring{\ml{കരടി} (karaṭi)}{കരടി (karaṭi)}]{\ml{കരടി} (karaṭi) — a bear}

\label{lesson:ML-C162-karati}



\begin{quote}
Before the new one: say the Malayalam for jasmine, then the Malayalam for a lotus.
\end{quote}

\subsection*{You'll want to know: \ml{കരടി}}
\textbf{\ml{കരടി}} — \emph{karaṭi} — \textquotedblleft{}a bear\textquotedblright{}.

\begin{cousinweb}
\textbf{\ml{കരടി}} (\emph{karaṭi}) is native Dravidian, one word with Tamil \emph{karaṭi}, Kannada \emph{karaḍi} and Telugu \emph{karaṭi} (DEDR 1263).
\end{cousinweb}

\begin{grammarlens}[title={What you've built}]
One more word for this chapter.
\end{grammarlens}

\subsection*{Your turn}
\begin{itemize}
\raggedright
  \item \emph{Say it:} \emph{karaṭi}
  \item \emph{Say it:} \emph{karaṭi}, once more
  \item \emph{Say it:} say \emph{karaṭi}
  \item \emph{Recall:} say the Malayalam for jasmine, then the Malayalam for a lotus, then say \emph{karaṭi} again
\end{itemize}

\begin{tcolorbox}[breakable,colback=teal!4,colframe=teal!35!black,title={Before you move on}]
What does \ml{കരടി} mean? (\textquotedblleft{}A bear\textquotedblright{}.) Say it once more.
\end{tcolorbox}
\section[\texorpdfstring{\ml{കുറുക്കൻ} (kuṟukkan)}{കുറുക്കൻ (kuṟukkan)}]{\ml{കുറുക്കൻ} (kuṟukkan) — a fox}

\label{lesson:ML-C162-kurukkan}



\begin{quote}
Before the new one: say the Malayalam for a lotus, then the Malayalam for a bear.

Say \textquotedblleft{}It's raining.\textquotedblright{} (\emph{Maḻa peyyunnu} — \textquotedblleft{}rain is falling.\textquotedblright{})
\end{quote}

\subsection*{You'll want to know: \ml{കുറുക്കൻ}}
\textbf{\ml{കുറുക്കൻ}} — \emph{kuṟukkan} — \textquotedblleft{}a fox\textquotedblright{}.

\begin{cousinweb}
\textbf{\ml{കുറുക്കൻ}} (\emph{kuṟukkan}) is native Dravidian, built on \emph{kuṟu}, \textquotedblleft{}short, little\textquotedblright{}: DEDR 1851 enters Malayalam \emph{kuṟukkan}, \textquotedblleft{}jackal\textquotedblright{}, under the root it shares with Tamil \emph{kuṟu}, \textquotedblleft{}short\textquotedblright{}.
\end{cousinweb}

\begin{grammarlens}[title={What you've built}]
One more word for this chapter.
\end{grammarlens}

\subsection*{Your turn}
\begin{itemize}
\raggedright
  \item \emph{Say it:} \emph{kuṟukkan}
  \item \emph{Say it:} \emph{kuṟukkan}, once more
  \item \emph{Say it:} say \emph{kuṟukkan}
  \item \emph{Recall:} say the Malayalam for a lotus, then the Malayalam for a bear, then say \emph{kuṟukkan} again
\end{itemize}

\begin{tcolorbox}[breakable,colback=teal!4,colframe=teal!35!black,title={Before you move on}]
What does \ml{കുറുക്കൻ} mean? (\textquotedblleft{}A fox\textquotedblright{}.) Say it once more.
\end{tcolorbox}
\section[\texorpdfstring{\ml{ഒട്ടകം} (oṭṭakaṁ)}{ഒട്ടകം (oṭṭakaṁ)}]{\ml{ഒട്ടകം} (oṭṭakaṁ) — a camel}

\label{lesson:ML-C162-ottakam}



\begin{quote}
Before the new one: say the Malayalam for a bear, then the Malayalam for a fox.
\end{quote}

\subsection*{You'll want to know: \ml{ഒട്ടകം}}
\textbf{\ml{ഒട്ടകം}} — \emph{oṭṭakaṁ} — \textquotedblleft{}a camel\textquotedblright{}.

\begin{cousinweb}
\textbf{\ml{ഒട്ടകം}} (\emph{oṭṭakaṁ}) goes back to Sanskrit \emph{uṣṭra}, \textquotedblleft{}camel\textquotedblright{} (CDIAL 2387), through a Prakrit form with \emph{ṭṭ}, with Malayalam's noun ending.
\end{cousinweb}

\begin{grammarlens}[title={What you've built}]
One more word for this chapter.
\end{grammarlens}

\subsection*{Your turn}
\begin{itemize}
\raggedright
  \item \emph{Say it:} \emph{oṭṭakaṁ}
  \item \emph{Say it:} \emph{oṭṭakaṁ}, once more
  \item \emph{Say it:} say \emph{oṭṭakaṁ}
  \item \emph{Recall:} say the Malayalam for a bear, then the Malayalam for a fox, then say \emph{oṭṭakaṁ} again
\end{itemize}

\begin{tcolorbox}[breakable,colback=teal!4,colframe=teal!35!black,title={Before you move on}]
What does \ml{ഒട്ടകം} mean? (\textquotedblleft{}A camel\textquotedblright{}.) Say it once more.
\end{tcolorbox}
\section[\texorpdfstring{\ml{താറാവ്} (tāṟāvŭ)}{താറാവ് (tāṟāvŭ)}]{\ml{താറാവ്} (tāṟāvŭ) — a duck}

\label{lesson:ML-C162-taravu}



\begin{quote}
Before the new one: say the Malayalam for a fox, then the Malayalam for a camel.
\end{quote}

\subsection*{You'll want to know: \ml{താറാവ്}}
\textbf{\ml{താറാവ്}} — \emph{tāṟāvŭ} — \textquotedblleft{}a duck\textquotedblright{}.

\begin{cousinweb}
\textbf{\ml{താറാവ്}} (\emph{tāṟāvŭ}) is native Dravidian, one word with Tamil \emph{tārā} (DEDR 3169).
\end{cousinweb}

\begin{grammarlens}[title={What you've built}]
One more word for this chapter.
\end{grammarlens}

\subsection*{Your turn}
\begin{itemize}
\raggedright
  \item \emph{Say it:} \emph{tāṟāvŭ}
  \item \emph{Say it:} \emph{tāṟāvŭ}, once more
  \item \emph{Say it:} say \emph{tāṟāvŭ}
  \item \emph{Recall:} say the Malayalam for a fox, then the Malayalam for a camel, then say \emph{tāṟāvŭ} again
\end{itemize}

\begin{tcolorbox}[breakable,colback=teal!4,colframe=teal!35!black,title={Before you move on}]
What does \ml{താറാവ്} mean? (\textquotedblleft{}A duck\textquotedblright{}.) Say it once more.
\end{tcolorbox}
\section[\texorpdfstring{\ml{മൂങ്ങ} (mūṅṅa)}{മൂങ്ങ (mūṅṅa)}]{\ml{മൂങ്ങ} (mūṅṅa) — an owl}

\label{lesson:ML-C162-munna}



\begin{quote}
Before the new one: say the Malayalam for a camel, then the Malayalam for a duck.

\emph{Mañña}, \textquotedblleft{}yellow,\textquotedblright{} and \emph{maññaḷ}, \textquotedblleft{}turmeric\textquotedblright{}: one root or two? (One — a native pair.)
\end{quote}

\subsection*{You'll want to know: \ml{മൂങ്ങ}}
\textbf{\ml{മൂങ്ങ}} — \emph{mūṅṅa} — \textquotedblleft{}an owl\textquotedblright{}.

\begin{cousinweb}
\textbf{\ml{മൂങ്ങ}} (\emph{mūṅṅa}) is of uncertain origin; DEDR does not list it.
\end{cousinweb}

\begin{grammarlens}[title={What you've built}]
One more word for this chapter.
\end{grammarlens}

\subsection*{Your turn}
\begin{itemize}
\raggedright
  \item \emph{Say it:} \emph{mūṅṅa}
  \item \emph{Say it:} \emph{mūṅṅa}, once more
  \item \emph{Say it:} say \emph{mūṅṅa}
  \item \emph{Recall:} say the Malayalam for a camel, then the Malayalam for a duck, then say \emph{mūṅṅa} again
\end{itemize}

\begin{tcolorbox}[breakable,colback=teal!4,colframe=teal!35!black,title={Before you move on}]
What does \ml{മൂങ്ങ} mean? (\textquotedblleft{}An owl\textquotedblright{}.) Say it once more.
\end{tcolorbox}
